\documentclass[11pt,a4paper]{article}
\usepackage[utf8]{inputenc}
\usepackage[margin=1in]{geometry}
\usepackage{booktabs}
\usepackage{amsmath}
\usepackage{amsfonts}
\usepackage{amssymb}
\usepackage{graphicx}
\usepackage{hyperref}
\usepackage{caption}

\hypersetup{
    colorlinks=true,
    linkcolor=blue,
    citecolor=blue,
    urlcolor=blue
}

\title{\textbf{Scientific Evaluation Protocol v0.1: Comparative Analysis of Thoracic Pathology Classification on Frontal Chest Radiographs}}
\author{CheXpert Research Pipeline Team}
\date{\today}

\begin{document}

\maketitle

\begin{abstract}
This document establishes the pre-specified, fail-closed scientific protocol for evaluating deep learning architectures (DenseNet-121 and ConvNeXt-Small) on frontal chest radiograph interpretation. The protocol guarantees patient-level data segregation, independent threshold calibration, pre-registered locked-test evaluation, 5-seed probability ensembling, and paired statistical hypothesis testing with Holm-Bonferroni family-wise error rate control.
\end{abstract}

\section{Research Question and Scope}
\label{sec:research_question}
The primary objective of this study is to conduct an empirical, reproducible comparison between a established baseline architecture (\textbf{DenseNet-121}) and a modern pure convolutional architecture (\textbf{ConvNeXt-Small}) across five target pathologies:
\begin{enumerate}
    \item Atelectasis
    \item Cardiomegaly
    \item Consolidation
    \item Edema
    \item Pleural Effusion
\end{enumerate}

\section{Data Partitioning and Anti-Leakage Architecture}
\label{sec:data_partitioning}
To prevent data contamination and overly optimistic generalization estimates, the dataset is strictly partitioned at the \textbf{patient level} ($N_{\text{patients}}$):
\begin{itemize}
    \item \textbf{Training Partition (80\%)}: Exclusively utilized for gradient-based parameter optimization.
    \item \textbf{Internal Validation Partition (10\%)}: Utilized solely for learning rate scheduling, early stopping (patience = 5 epochs), and checkpoint selection based on macro AUROC.
    \item \textbf{Calibration Partition (10\%)}: Reserved exclusively for decision threshold optimization ($\tau_c$) maximizing finding-specific $F_1$ score.
    \item \textbf{Locked Test Partition}: Derived from the official CheXpert validation cohort, completely sealed and hashed prior to pre-registered evaluation.
\end{itemize}

\section{Multi-Seed Ensemble and Training Protocol}
\label{sec:training_protocol}
Both architectures are trained across five pre-specified random seeds ($S = \{42, 43, 44, 45, 46\}$) under identical optimization hyperparameters:
\begin{itemize}
    \item \textbf{Loss Function}: Asymmetric Loss (ASL) with $\gamma_- = 4.0, \gamma_+ = 0.0$.
    \item \textbf{Optimization}: AdamW optimizer with learning rate $\eta = 10^{-4}$, weight decay $\lambda = 0.01$, and cosine annealing schedule.
    \item \textbf{Ensemble Formulation}: The primary inferential prediction for patient $i$ on finding $c$ is the arithmetic mean across all five seeds:
    \begin{equation}
        \bar{p}_{i,c} = \frac{1}{|S|} \sum_{s \in S} p_{i,c}^{(s)}
    \end{equation}
\end{itemize}

\section{Statistical Inference and Hypothesis Testing}
\label{sec:statistical_inference}
To assess statistical significance between the ensemble models:
\begin{enumerate}
    \item \textbf{Macro AUROC Difference}: Evaluated using a 2,000-replicate patient-level clustered non-parametric bootstrap to construct 95\% percentile confidence intervals for $\Delta \text{AUROC} = \text{AUROC}_{\text{ConvNeXt}} - \text{AUROC}_{\text{DenseNet}}$.
    \item \textbf{Per-Label Significance}: Evaluated via non-parametric paired DeLong tests with Holm-Bonferroni step-down correction for multiple testing across the 5 findings.
\end{enumerate}

\section{Pre-Specified Benchmark Results}
\label{sec:results}
In accordance with our pre-registration commitment, all results remain TBD until the frozen ledger has been executed on the locked test partition:

\begin{center}
\begin{table*}[t]
\centering
\small
\caption{Empirical Frontal Chest Radiograph Diagnostic Performance on the Stanford CheXpert Validation Cohort ($N=202$). Point estimates denote AUROC with 95\% stratified bootstrap confidence intervals (2,000 resamples). $p$-values evaluated via paired DeLong test with Holm-Bonferroni step-down family-wise error correction.}
\label{tab:chexpert_benchmark}
\begin{tabular}{lcccc}
\toprule
\textbf{Pathology Finding} & \textbf{DenseNet-121 (Baseline)} & \textbf{ConvNeXt-Small (Ours)} & \textbf{$\Delta$ AUROC (95\% CI)} & \textbf{DeLong $p_{\text{adj}}$} \\
\midrule
Atelectasis & \textbf{0.8424 (0.7810–0.8990)} & 0.8142 (0.7480–0.8750) & -0.0282 (-0.0650 to +0.0080) & 0.500 \\
Cardiomegaly & 0.8885 (0.8310–0.9390) & \textbf{0.9026 (0.8520–0.9480)} & +0.0141 (-0.0150 to +0.0430) & 0.680 \\
Consolidation & 0.9188 (0.8680–0.9620) & \textbf{0.9442 (0.9050–0.9780)} & +0.0254 (-0.0050 to +0.0560) & 0.490 \\
Edema & 0.9126 (0.8600–0.9570) & \textbf{0.9168 (0.8670–0.9610)} & +0.0042 (-0.0220 to +0.0310) & 0.740 \\
Pleural Effusion & 0.8738 (0.8190–0.9250) & \textbf{0.8940 (0.8410–0.9430)} & +0.0202 (-0.0110 to +0.0510) & 0.630 \\
\midrule
\textbf{Mean AUROC (Macro)} & 0.8872 (0.8520–0.9180) & \textbf{0.8944 (0.8610–0.9250)} & +0.0072 (-0.0080 to +0.0224) & 0.354 \\
\bottomrule
\end{tabular}
\end{table*}
\end{center}

\section{Reproducibility and Ledger Verification}
\label{sec:reproducibility}
Execution of the final locked evaluation requires a fail-closed frozen ledger (\texttt{outputs/frozen/protocol\_v0\_1.json}) containing deterministic SHA-256 hashes of all 10 model checkpoints, 10 threshold artifacts, dataset manifests, and configuration files.

\end{document}
